\chapter{A Generic Local Nonlinear Problem Framework for Constitutive Updates}

\section{Motivation}

The previous refactor split the classical small-strain plastic update into
three compile-time axes:
\begin{itemize}
  \item the yield function,
  \item the consistency residual and its Jacobian,
  \item the return algorithm.
\end{itemize}

That split is already useful, but it still leaves one structural assumption in
place: the local nonlinear problem is scalar.  In the current $J_2$ path the
unknown is only the plastic-multiplier increment $\Delta\gamma$, and the local
solve is
\[
  r(\Delta\gamma) = 0.
\]

This is too narrow for the next generation of constitutive models that the
library aims to support:
\begin{itemize}
  \item finite-strain plasticity based on multiplicative decomposition,
  \item coupled damage-plasticity,
  \item pressure-sensitive plasticity with additional local variables,
  \item local fracture regularization variables,
  \item rate-dependent or mixed local constitutive updates.
\end{itemize}

In those models the local unknown is generally a vector
\[
  \mathbf{x} \in \mathbb{R}^m,
\]
and the constitutive update is written as a local nonlinear system
\[
  \mathbf{R}(\mathbf{x}) = \mathbf{0},
\]
with a consistent local Jacobian
\[
  \mathbf{K}(\mathbf{x}) =
  \frac{\partial \mathbf{R}}{\partial \mathbf{x}}.
\]

The right architectural move is therefore to promote the \emph{local nonlinear
problem itself} to a first-class compile-time abstraction.

\section{Architectural Split}

This phase introduces three layers.

\subsection{Local Problem Contract}

The new header
\texttt{src/materials/local\_problem/LocalNonlinearProblem.hh}
defines a generic concept:
\begin{itemize}
  \item \texttt{UnknownT}: the local unknown type,
  \item \texttt{ResidualT}: the local residual type,
  \item \texttt{JacobianT}: the local Jacobian type,
  \item \texttt{initial\_guess(relation, context)},
  \item \texttt{residual(relation, context, x)},
  \item \texttt{jacobian(relation, context, x)},
  \item \texttt{residual\_norm(r)}.
\end{itemize}

The design intentionally keeps the \texttt{context} type open.  For the current
small-strain plasticity path it stores the trial state, the algorithmic state,
the flow direction, and the trial overstress.  For future finite-strain models
the same slot can carry
\[
  \mathbf{F}_{n+1}, \quad \mathbf{C}_{n+1}, \quad \mathbf{b}_{n+1}, \quad
  \Delta t, \quad \alpha_n,
\]
or any other local constitutive data required by the update.

\subsection{Generic Newton Solver}

The header
\texttt{src/materials/local\_problem/NewtonLocalSolver.hh}
implements a reusable Newton iteration:
\[
  \mathbf{K}(\mathbf{x}^{(i)}) \, \Delta \mathbf{x}^{(i)}
  = -\mathbf{R}(\mathbf{x}^{(i)}),
\]
\[
  \mathbf{x}^{(i+1)} = \mathbf{x}^{(i)} + \Delta \mathbf{x}^{(i)}.
\]

The implementation is generic over the local unknown type:
\begin{itemize}
  \item if the unknown is scalar, the update reduces to
        $-\!R/K$,
  \item if the unknown is a fixed-size Eigen vector, the update is solved as a
        small dense linear system.
\end{itemize}

This matters for HPC-oriented design because the type of the local unknown
remains fixed at compile time, and the solver does not introduce virtual
dispatch into the constitutive hot path.

\subsection{Scalar Consistency as an Adapter}

The header
\texttt{src/materials/constitutive\_models/non\_lineal/plasticity/ScalarConsistencyProblem.hh}
turns the current scalar consistency update into a \texttt{LocalNonlinearProblem}.

Thus, the current plasticity route is reinterpreted as the special case
\[
  x = \Delta\gamma, \qquad
  R(x) = r(\Delta\gamma), \qquad
  K(x) = \frac{\mathrm{d} r}{\mathrm{d}\Delta\gamma}.
\]

This is the important architectural result: the present implementation remains
valid, but it now lives inside a more general framework that no longer assumes
scalar local updates.

\section{Return Mapping as a Specialization of the General Framework}

The new
\texttt{LocalNonlinearReturnAlgorithm<LocalProblem, LocalSolver>}
in
\texttt{src/materials/constitutive\_models/non\_lineal/plasticity/ReturnAlgorithm.hh}
acts as the generic bridge between:
\begin{itemize}
  \item a constitutive relation that can build a local update context,
  \item a local problem definition,
  \item a local nonlinear solver.
\end{itemize}

The old user-facing
\texttt{NewtonConsistencyReturnAlgorithm<Residual, Jacobian, MaxIterations>}
is kept as a compatibility wrapper.  Internally, it now instantiates
\[
  \texttt{ScalarConsistencyProblem<Residual, Jacobian>}
\]
and
\[
  \texttt{NewtonLocalSolver<MaxIterations>}.
\]

This gives two immediate benefits:
\begin{itemize}
  \item backward compatibility for the existing small-strain $J_2$ path,
  \item a direct path to future vector-valued local problems.
\end{itemize}

\section{Why This Matters for Large Displacements and Finite Strains}

Finite-strain inelastic constitutive updates are typically not expressible as a
single scalar consistency correction.  In multiplicative elastoplasticity, for
example, one often solves for combinations of:
\begin{itemize}
  \item a plastic multiplier increment,
  \item an updated plastic metric or elastic left Cauchy--Green tensor,
  \item hardening variables,
  \item damage or fracture-driving variables.
\end{itemize}

The local problem then takes the block form
\[
  \mathbf{R}(\mathbf{x}) =
  \begin{bmatrix}
    R_1(\mathbf{x}) \\
    R_2(\mathbf{x}) \\
    \vdots \\
    R_m(\mathbf{x})
  \end{bmatrix}
  = \mathbf{0},
\]
with a block Jacobian
\[
  \mathbf{K}(\mathbf{x}) =
  \begin{bmatrix}
    \partial R_1/\partial x_1 & \cdots & \partial R_1/\partial x_m \\
    \vdots & \ddots & \vdots \\
    \partial R_m/\partial x_1 & \cdots & \partial R_m/\partial x_m
  \end{bmatrix}.
\]

The current refactor does not yet implement a full finite-strain plastic law,
but it removes a key architectural obstacle: the local constitutive solver no
longer assumes that the unknown is scalar.

\section{Verification}

This phase adds two important regression checks:
\begin{itemize}
  \item the generic
        \texttt{LocalNonlinearReturnAlgorithm<>}
        reproduces the current small-strain radial-return response and tangent;
  \item the generic Newton local solver converges on a fixed-size vector
        two-equation problem independent of plasticity.
\end{itemize}

The second test is intentionally simple.  Its value is architectural: it proves
that the new framework already supports vector local problems without any extra
branching in the solver.

\section{Strategic Outcome}

At this point the local constitutive stack is factored as:
\begin{itemize}
  \item \texttt{YieldFunction}: admissibility scalar,
  \item \texttt{ConsistencyResidual / Jacobian}: local equation ingredients,
  \item \texttt{ReturnAlgorithm}: correction strategy,
  \item \texttt{LocalNonlinearProblem}: abstract local system definition,
  \item \texttt{NewtonLocalSolver}: reusable nonlinear solver.
\end{itemize}

This is the right level of generality for the next stages:
\begin{itemize}
  \item finite-strain elastoplasticity with vector local unknowns,
  \item damage-plasticity with coupled internal variables,
  \item fracture-oriented local updates regularized at constitutive level,
  \item objective-rate or multiplicative algorithms integrated through the same
        erased \texttt{Material<>} boundary.
\end{itemize}

\section{References}

\begin{thebibliography}{9}

\bibitem{simo1998localproblem}
J.~C. Simo and T.~J.~R. Hughes,
\emph{Computational Inelasticity},
Springer, 1998.

\bibitem{desouza2008localproblem}
E.~A. de Souza Neto, D.~Peri\'{c}, and D.~R.~J. Owen,
\emph{Computational Methods for Plasticity: Theory and Applications},
Wiley, 2008.

\bibitem{miehe1996}
C.~Miehe,
``Numerical computation of algorithmic (consistent) tangent moduli in large-strain computational inelasticity,''
\emph{Computer Methods in Applied Mechanics and Engineering},
134(3--4):223--240, 1996.

\end{thebibliography}
